\noindent\textbf{(a) Resultado.} $T^{\prime}=-k(T-20)$ con $k>0$; $k$ se mide en $\mathrm{min}^{-1}$.

\noindent\textit{Desarrollo.} Si la capacidad calorífica total es $C_H>0$ y la potencia de salida es $H(T-20)$, con $H>0$, el balance es $C_HT^{\prime}=-H(T-20)$. Con $k=H/C_H>0$ resulta el modelo indicado. Si $T>20$, la derivada es negativa; si $T<20$, el signo se invierte y el entorno calienta el cuerpo. En la convención $T^{\prime}=\alpha(T-20)$, la constante sería $\alpha=-k<0$.

\medskip

\noindent\textbf{(b) Resultado.} $T(60)=22.8125\,{}^\circ\mathrm C$; el descenso desde el minuto 30 es $12.1875\,{}^\circ\mathrm C$.

\noindent\textit{Desarrollo.} Separar da $T=20+80e^{-kt}$. El dato de 30 minutos impone $e^{-30k}=15/80=3/16$, luego $k=\ln(16/3)/30\ \mathrm{min}^{-1}$. Por tanto $T(60)=20+80(3/16)^2=20+45/16=22.8125$. El descenso en ese intervalo es $35-T(60)=195/16=12.1875$. El dato inicial y el de calibración se recuperan al sustituir. La temperatura desciende hacia 20 sin cruzarlo; el segundo descenso es menor que el primero, de acuerdo con la menor diferencia térmica.
